\documentclass[11pt]{article}
\input{packages.tex}
\usepackage{ulem}
\DeclareMathOperator{\REG}{REG}


% Настройка колонтитулов
\pagestyle{fancy}
\fancyhf{} % Очистить все поля
\fancyhead[L]{БУНТ}
\fancyhead[R]{\href{https://t.me/sol_m_07}{Соля}}
\fancyfoot[C]{\thepage}               % Справа внизу

\setlength{\headheight}{13.6pt}

\hypersetup{
    colorlinks=true,
    linkcolor=blue,
    citecolor=blue,
    filecolor=blue,
    urlcolor=blue,
}

\begin{document}

\begin{titlepage}
    \centering
    \vspace*{\fill}  % заполняет пространство сверху

    {\Huge Лекция 6. Формы записи регулярных языков
    \par}

    \vspace*{\fill}  % заполняет пространство снизу
\end{titlepage}

\setcounter{page}{1}

\section{Блок повторения}

Так как на лекции сначала было повторение материала прошлого конспекта, 
то автор советует его перечитать перед погружением в доказательство.

\section{Формы записи регулярных языков}

\begin{theorem}
    Пусть $\Sigma$ -- алфавит. Тогда имеем $2^{\Sigma^{\ast}}$ языков, 
    заданных над ним. $\REG \subseteq 2^{\Sigma^{\ast}}$. Каждый 
    регулярный язык может быть записан через
    \begin{enumerate}
        \item регулярные выражения
        \item ДКА
        \item НКА
        \item грамматику
    \end{enumerate}
    Причём из одной формы записи мы всегда можем получить другую. 
    Докажем это в следующих пунктах.
\end{theorem}

\subsection{От регулярного выражения к НКА}

Введём правила для НКА:
\begin{enumerate}
    \item НКА имеет ровно одно принимающее состояние
    \item В начальное состояние не ведёт ни один переход
    \item Из принимающего состояния нет ни одного перехода
\end{enumerate}

\begin{lemma}
    Любой НКА можно подогнать под наши правила.
\end{lemma}

\begin{proof}
    Пойдём по порядку. Если у нас изначально одно принимающее 
    состояние, ничего не делаем. Если же принимающих состояний 
    несколько, то мы делаем их "обычными" \ и строим из них 
    $\varepsilon \text{-переходы}$ в единое принимающее состояние.
    \begin{figure}[h]
        \centering
        \includegraphics[width=0.9\linewidth]{pictures/НКА1.png}
    \end{figure} \leavevmode
    \\ \\
    Заметим, что третье правило мы возможно уже выполнили. Если ещё 
    не выполнили, то по такому же принципу образуем новое принимающее 
    состояние и $\varepsilon \text{-переход}$ к нему от текущего 
    (само текущее делаем "обычным"). \\
    \begin{figure}[h]
        \centering
        \includegraphics[width=0.9\linewidth]{pictures/НКА2.png}
    \end{figure} \leavevmode
    \\
    Осталось вернуть долг второму правилу. Здесь в случае невыполнения 
    мы образуем новое начальное состояние и из него строим 
    $\varepsilon \text{-переходы}$ в текущие начальные состояния. \\
    \begin{figure}[h]
        \centering
        \includegraphics[width=1\linewidth]{pictures/НКА3.png}
    \end{figure}
\end{proof} \leavevmode
\\
Будем схематично изображать НКА языка $X$ как \\
\begin{figure}[h]
    \centering
    \includegraphics[width=0.7\linewidth]{pictures/НКА.png}
\end{figure} \leavevmode
\\ \\ \\ \\ \\ \\
Тогда
\begin{figure}[h]
    \centering
    \includegraphics[width=0.5\linewidth]{pictures/Конкатенация.png}
    \caption{Конкатенация $X \cdot Y$}
\end{figure}
\begin{figure}[h]
    \centering
    \includegraphics[width=0.55\linewidth]{pictures/Объединение.png}
    \caption{Объединение $X + Y$}
\end{figure}
\begin{figure}[h]
    \centering
    \includegraphics[width=0.55\linewidth]{pictures/Итерация.png}
    \caption{Итерация $X^*$}
\end{figure} \leavevmode 
\\ \\
Алгоритм построения будет заключаться в построении дерева по 
приоритету операций, а затем постепенного объединения результатов 
снизу вверх с помощью схем, представленных выше. 
\\
\begin{example}
    $\Sigma = \{a, b\}$. Рассмотрим регулярное выражение 
    $((ab + b^*)aa)^*$. Составим дерево по приоритету выполнения 
    операций и построим НКА снизу вверх:
    \begin{figure}[H]
        \centering
        \includegraphics[width=0.4\linewidth]{pictures/Дерево.png}
    \end{figure}
    \begin{figure}[H]
        \centering
        \includegraphics[width=0.8\linewidth]{pictures/Построение НКА.png}
    \end{figure}
    \begin{figure}[H]
        \centering
        \includegraphics[width=1\linewidth]{pictures/Завершённый НКА.png}
    \end{figure}
\end{example} \leavevmode
\\
\begin{remark}
    Нумерация состояний пригодится нам в построении ДКА. Все 
    дальнейшие примеры будут использовать эту схему.
\end{remark}

\begin{remark}
    Очень важно использовать готовые схемы для конкатенации, 
    объединения и итерации. Они будут приниматься без доказательств. 
    Собственные схемы необходимо будет доказать (обычно получается 
    лажа).
\end{remark}

\subsection{От НКА к ДКА}

Вход: НКА($\Sigma, Q, Q_f, \delta, q_0$) \\
Выход: ДКА($\Sigma, Q', Q_f', \delta', q_0'$) \\
Состояния ДКА будут соответствовать подмножествам состояний 
исходного НКА.

\begin{definition}
    $\varepsilon \text{-замыканием}$ множества состояний $Q$ будем 
    называть множество состояний, достижимых из $Q$ по $\varepsilon 
    \text{-переходам}$, и обозначать $C_{\varepsilon}(Q)$
\end{definition}

\begin{example}
    Рассмотрим $\varepsilon \text{-замыкания}$ некоторых множеств 
    состояний:
    \begin{itemize}
        \item $C_{\varepsilon}(\{1\}) = \{1, 2, 6, 7, 9, 5\}$
        \item $C_{\varepsilon}(\{0, 1\}) = \{1, 2, 6, 7, 9, 5, 0, 12\}$
    \end{itemize}
\end{example}

\begin{definition}
    Определим множество состояний, в которые можно перейти из 
    $R \subseteq Q$ при считывании одного $a \in \Sigma$:
    \[move(R, a) := \{q' \ | \ \exists q \in R : \delta (a, q) = q'\}\]
\end{definition}

\begin{example}
    $move(\{2, 8\}, a) = \{3\}$
\end{example} \leavevmode
\\
Пора поведать алгоритм построения ДКА. Сначала $Q'$ и $\delta'$ пусты. 
\begin{enumerate}
    \item Определим $q_0' = C_{\varepsilon}(q_0)$
    \item Добавляем $q_0'$ в $Q'$ как непомеченное состояние
    \item Пока в $Q'$ есть непомеченное $R$:
    \begin{itemize}
        \item Пометить $R$
        \item Для каждого символа $a \in \Sigma$
        \begin{itemize}
            \item $S = C_{\varepsilon}(move(R, a))$
            \item Если $S \neq \varnothing$:
            \begin{itemize}
                \item Если $S \notin Q'$:
                \begin{itemize}
                    \item Добавить $S$ в $Q'$ как непомеченное состояние
                \end{itemize}
                \item Определить $\delta'(a, R) = S$
            \end{itemize}
        \end{itemize}
    \end{itemize}
    \item Определить $Q_f' = \{S \ | \ S \in Q', \ S \cap Q_f \neq \varnothing\}$ 
\end{enumerate} \leavevmode
\\
\begin{example}
    $q_0' = C_{\varepsilon}(0) = \{0, 12, 1, 2, 6, 7, 9, 5\} \longrightarrow Q' = \{q_0'\}$ \\ \\
    $S = C_{\varepsilon}(move(q_0', a)) = C_{\varepsilon}(\{3, 10\}) = \{3, 10\} = q_1' \longrightarrow Q' = \{q_0', q_1'\}$ \\ \\
    $S = C_{\varepsilon}(move(q_0', b)) = C_{\varepsilon}(\{8\}) = \{8, 9, 7, 5\} = q_2' \longrightarrow Q' = \{q_0', q_1', q_2'\}$ \\ \\
    $S = C_{\varepsilon}(move(q_1', a)) = C_{\varepsilon}(\{11\}) = \{11, 12, 1, 2, 6, 7, 9, 5\} = q_3' \longrightarrow Q' = \{q_0', q_1', q_2', q_3'\}$ \\ \\
    $S = C_{\varepsilon}(move(q_1', b)) = C_{\varepsilon}(\{4\}) = \{5\} = q_4' \longrightarrow Q' = \{q_0', q_1', q_2', q_3', q_4'\}$ \\ \\
    $S = C_{\varepsilon}(move(q_2', a)) = C_{\varepsilon}(\{10\}) = \{10\} = q_5' \longrightarrow Q' = \{q_0', q_1', q_2', q_3', q_4', q_5'\}$ \\ \\
    $S = C_{\varepsilon}(move(q_2', b)) = C_{\varepsilon}(\{8\}) = \{8, 9, 5, 7\} = q_2'$ \\ \\
    $S = C_{\varepsilon}(move(q_3', a)) = C_{\varepsilon}(\{3, 10\}) = \{3, 10\} = q_1'$ \\ \\ 
    $S = C_{\varepsilon}(move(q_3', b)) = C_{\varepsilon}(\{8\}) = \{8, 9, 7, 5\} = q_2'$ \\ \\
    $S = C_{\varepsilon}(move(q_4', a)) = C_{\varepsilon}(\{10\}) = \{10\} = q_5'$ \\ \\
    $S = C_{\varepsilon}(move(q_4', b)) = C_{\varepsilon}(\varnothing) = \varnothing$ \\ \\
    $S = C_{\varepsilon}(move(q_5', a)) = C_{\varepsilon}(\{11\}) = \{11, 12, 1, 2, 6, 7, 9, 5\} = q_3'$ \\ \\
    $S = C_{\varepsilon}(move(q_5', b)) = C_{\varepsilon}(\varnothing) = \varnothing$ \\ \\ \\
    $Q' = \{q_0', q_1', q_2', q_3', q_4', q_5'\}, \quad Q_f' = \{q_0', q_3'\}$ 
    \begin{figure}[H]
        \centering
        \includegraphics[width=0.5\linewidth]{pictures/ДКА.png}
    \end{figure}
\end{example}

\section{Решение системы уравнений с регулярными коэффициентами}

\[\begin{cases}
    X_1 = \alpha_{11} X_1 + \dots + \alpha_{1n} X_n + \beta_1 \\
    \vdots \\
    X_n = \alpha_{n1} X_1 + \dots + \alpha_{nn} X_n + \beta_n
\end{cases}\]

Здесь $X_i$ -- переменные множества, а $\alpha_{ij}$ и $\beta_i$ -- 
это регулярные языки, задаваемые регулярными выражениями. Решением 
будем называть набор языков $A_1, \dots, A_n$ т.ч. для любого набора 
языков $B_1, \dots, B_n$, удовлетворяющего системе, выполнено $A_i 
\subseteq B_i$ (решение является минимальным по включению набором).
\\ \\ 
Разберём случаи при решении системы:

\paragraph{Случай 1}

Если в системе есть строка вида 
\[X_i = a_{ii} X_i + \gamma, \ a_{ii} \neq \varnothing,\]
то заменим всюду $X_i$ на $a_{ii}^* \gamma$, а саму строку удалим.

\begin{example}
    \[X = aX + b \implies X = a^* b\]
\end{example}

\paragraph{Случай 2}

Если в системе нет строк, подходящих под первый случай, то переменная 
$X_n$ не встречается в правой части $n \text{-ой}$ строки. Заменим 
всюду $X_n$ на правую часть этой строки. Саму строку удалим.

\section{Правые языки}

\begin{definition}
    По заданному НКА для каждого $q$ определим правый язык:
    \[R_q := \{x \ | \ q \xrightarrow{x} q_f, \ q_f \in Q_f\}\]
    То есть все слова, считывая которые, можно попасть из состояния 
    $q$ в одно из финальных состояний. 
\end{definition} \leavevmode
\\
Именно правые языки мы будем использовать в качестве переменных 
при составлении системы. \\ \\ 
Для каждого состояния $q_i$ и каждого сивола $a$: если из состояния 
$q_i$ есть переход по символу $a$ в $q_j$, то к правой части строки 
$R_{q_i}$ добавим $+aR_{q_j}$. К правым частям финальных состояний 
добавим $+ \varepsilon$.

\begin{example}
    \begin{figure}[H]
        \centering
        \includegraphics[width=0.5\linewidth]{pictures/Правые языки.png}
    \end{figure}
    \[\begin{cases}
        R_0 = aR_1 \\
        R_1 = bR_0 + \varepsilon
    \end{cases}\]
    Согласно алгоритму: $R_0 = a(bR_0 + \varepsilon) = abR_0 + a 
    \implies R_0 = (ab)^* a$
\end{example}

\section{Мысли вслух}

\begin{figure}[H]
    \centering
    
\includegraphics[width=0.9\linewidth]{pictures/Ёжик.jpg}
\end{figure}

\end{document}